% ============================================================
% 04-bezier-matrix-straight-output-explained.tex
% Compile with:  pdflatex -shell-escape 04-bezier-matrix-straight-output-explained.tex
% Run twice for table of contents.
% ============================================================

\documentclass[12pt, a4paper]{article}

% ---- Encoding & fonts ----------------------------------------
\usepackage[T1]{fontenc}
\usepackage[utf8]{inputenc}
\usepackage{lmodern}

% ---- Page geometry -------------------------------------------
\usepackage[
  top=2.5cm, bottom=2.5cm,
  left=2.8cm, right=2.8cm
]{geometry}

% ---- Colours -------------------------------------------------
\usepackage[dvipsnames]{xcolor}
\definecolor{codebg}{RGB}{248,248,248}
\definecolor{rulecolor}{RGB}{180,180,180}
\definecolor{examplebg}{RGB}{240,248,255}
\definecolor{examplerule}{RGB}{100,149,237}
\definecolor{notebg}{RGB}{255,250,230}
\definecolor{noterule}{RGB}{210,180,50}

% ---- Minted --------------------------------------------------
\usepackage{minted}
\setminted[text]{
  bgcolor    = codebg,
  frame      = leftline,
  framerule  = 1.5pt,
  rulecolor  = rulecolor,
  fontsize   = \small,
  breaklines = true,
}

% ---- Math ----------------------------------------------------
\usepackage{amsmath}
\usepackage{amssymb}

% ---- Boxes ---------------------------------------------------
\usepackage{mdframed}

\newmdenv[
  backgroundcolor  = examplebg,
  linecolor        = examplerule, linewidth=1.5pt,
  leftline=true, rightline=false, topline=false, bottomline=false,
  innerleftmargin=10pt, innerrightmargin=8pt,
  innertopmargin=6pt, innerbottommargin=6pt,
  skipabove=6pt, skipbelow=4pt,
]{examplebox}

\newmdenv[
  backgroundcolor  = notebg,
  linecolor        = noterule, linewidth=1.5pt,
  leftline=true, rightline=false, topline=false, bottomline=false,
  innerleftmargin=10pt, innerrightmargin=8pt,
  innertopmargin=6pt, innerbottommargin=6pt,
  skipabove=6pt, skipbelow=4pt,
]{notebox}

% ---- Hyperlinks ----------------------------------------------
\usepackage[colorlinks=true, linkcolor=black, urlcolor=MidnightBlue]{hyperref}

% ---- Misc ----------------------------------------------------
\usepackage{booktabs}
\usepackage{needspace}
\usepackage{parskip}
\usepackage{titlesec}
\usepackage{fancyhdr}
\usepackage{datetime2}

% ---- Section formatting --------------------------------------
\titleformat{\section}
  {\large\bfseries\color{MidnightBlue}}
  {\thesection.}{0.6em}{}[\vspace{-4pt}\rule{\linewidth}{0.4pt}]
\titleformat{\subsection}
  {\normalsize\bfseries\color{BrickRed}}
  {\thesubsection.}{0.5em}{}

% ---- Header / footer -----------------------------------------
\pagestyle{fancy}
\fancyhf{}
\rhead{\textcolor{gray}{\small 04-bezier-matrix-straight-output-explained}}
\lhead{\textcolor{gray}{\small B\'ezier Matrix Straight --- Output Explanation}}
\cfoot{\thepage}
\renewcommand{\headrulewidth}{0.3pt}

% ---- Helpers -------------------------------------------------
\newcommand{\cd}[1]{\texttt{#1}}

% ---- Title ---------------------------------------------------
\title{%
  \vspace{1cm}
  {\LARGE\bfseries B\'ezier Curve --- Matrix Method (Straight)}\\[0.4em]
  {\large $P(u) = U \cdot M \cdot G$}\\[0.4em]
  {\large Explanation of the Program Output}\\[0.4em]
  {\normalsize Case 2: Quartic 3-D, 5 Control Points}\\[0.4em]
  {\normalsize\texttt{04-bezier-matrix-straight.jl}}
}
\author{E.R.\ Nurwijayadi}
\date{\DTMtoday}


% =============================================================
\begin{document}
% =============================================================

\maketitle
\thispagestyle{empty}

\begin{center}
\textit{Directed and curated by E.R.\ Nurwijayadi.
Code and explanations AI-generated.}
\end{center}

\begin{center}
\begin{minipage}{0.85\textwidth}
\small
This document explains the output of
\cd{04-bezier-matrix-straight.jl} when run with \cd{CASE = 2}.
The mathematical method is identical to script~03
($P(u) = U \cdot M \cdot G$, same control points, same matrices,
same results).
What changed is the \textbf{output format}: script~03 was
pedagogical --- it derived $M$ column by column, showed each
dot product explicitly, and printed five numbered steps.
Script~04 is straight --- it prints the final matrices and
result directly, with no intermediate working shown.
\end{minipage}
\end{center}

\vspace{1em}\hrule\vspace{1.5em}
\tableofcontents
\newpage


% =============================================================
\needspace{8\baselineskip}
\section{What Changed from Script 03}
% =============================================================

\textit{Same curve. Same matrices. Same numbers.
Just a cleaner, more direct presentation.}

Script~03 had five numbered steps, each showing the derivation
or computation in detail.
Script~04 has five \emph{sections} --- flat, unnumbered,
one per matrix or result --- printed without any working.

\begin{center}
\begin{tabular}{lll}
\toprule
Script~03 output block & Script~04 output block & Content \\
\midrule
Step~1: Matrix Formulation & \textit{(header only)} & Formula and shapes \\
Step~2: Derive $M$ & $M$ --- Basis Matrix & Just the matrix \\
Step~3: Compute $M \cdot G$ & $M \cdot G$ --- Polynomial Coefficient Matrix & Just the matrix \\
\textit{(inside Step~3)} & Resulting Polynomials & The polynomial strings \\
Step~4: Evaluate $U \cdot (M \cdot G)$ & \textit{(dropped)} & Per-$u$ dot products \\
Step~5: Summary table & Evaluation Table & Box-bordered table \\
\midrule
\textit{(no equivalent)} & $G$ --- Geometry Matrix & Control points as matrix \\
\bottomrule
\end{tabular}
\end{center}

\begin{notebox}
\textbf{One new block: $G$ is now shown explicitly.}
Script~03 listed the control points as a plain list in Step~1.
Script~04 prints them as the labelled geometry matrix $G$
--- the same data, formatted as a proper matrix table,
making the connection to $P(u) = U \cdot M \cdot G$ more direct.

\textbf{One dropped block: the per-$u$ dot product display.}
Script~03 Step~4 showed every multiplication
$U[k] \cdot (MG)[k]$ for each $u$ value.
Script~04 drops this entirely and goes straight to the
summary table.
The numbers are the same; you just no longer see how they
were reached.
\end{notebox}


% =============================================================
\needspace{8\baselineskip}
\section{Header}
% =============================================================

\begin{minted}{text}
==================================================================
  Bezier Curve - Matrix Method  [P(u) = U * M * G]
  Case 2  |  Quartic  |  3D  |  5 Control Points
==================================================================
\end{minted}

Identical to script~03.
The formula $P(u) = U \cdot M \cdot G$ is stated in the
title line, setting the context for the four matrices
that follow.


% =============================================================
\needspace{8\baselineskip}
\section{$G$ --- Geometry Matrix}
% =============================================================

\begin{minted}{text}
  G  - Geometry Matrix (Control Points)
  ----------------------------------------------------------------

      x  y  z
      -  -  -
  P0  0  0  1
  P1  0  4  1
  P2  4  0  1
  P3  4  4  1
  P4  5  4  1
\end{minted}

$G$ is the $5 \times 3$ control point matrix:

\[
  G = \begin{bmatrix}
    0 & 0 & 1 \\
    0 & 4 & 1 \\
    4 & 0 & 1 \\
    4 & 4 & 1 \\
    5 & 4 & 1
  \end{bmatrix}
\]

Row $i$ = control point $P_{i-1}$.
Column $j$ = axis ($x$, $y$, $z$).
This is the same data as the control point list in the
header of scripts~01--03, now presented as the matrix
it actually is in the computation $P(u) = U \cdot M \cdot G$.

All $z$-coordinates are 1, which means the curve lies
entirely in the plane $z = 1$. This will be confirmed
when $z(u) = 1$ appears in every result below.


% =============================================================
\needspace{8\baselineskip}
\section{$M$ --- Basis Matrix}
% =============================================================

\begin{minted}{text}
  M  - Basis Matrix (derived from Bernstein)
  ----------------------------------------------------------------

       P0   P1   P2  P3  P4
       --  ---  ---  --  --
  u^4   1   -4    6  -4   1
  u^3  -4   12  -12   4   0
  u^2   6  -12    6   0   0
  u^1  -4    4    0   0   0
  u^0   1    0    0   0   0
\end{minted}

$M$ is the $5 \times 5$ basis matrix derived from the
Bernstein polynomials.
It is the same matrix that script~03 derived column by
column in Step~2.
Here it is printed directly without derivation.

\[
  M = \begin{bmatrix}
    1  & -4 &  6 & -4 & 1 \\
   -4  & 12 & -12 & 4 & 0 \\
    6  & -12 &  6 &  0 & 0 \\
   -4  &  4 &  0 &  0 & 0 \\
    1  &  0 &  0 &  0 & 0
  \end{bmatrix}
\]

The column labels \cd{P0} through \cd{P4} indicate that
column $i$ holds the monomial coefficients of $B(4,i,u)$,
the $i$-th Bernstein basis polynomial.
The row labels indicate the power of $u$ each row represents,
from $u^4$ (top) to $u^0 = 1$ (bottom).

\begin{notebox}
\textbf{Reading $M$:}
Column~0 contains $[1, -4, 6, -4, 1]$ --- the coefficients
of $(1-u)^4 = 1 - 4u + 6u^2 - 4u^3 + u^4$,
reading from the $u^4$ row downward.
Column~4 contains $[1, 0, 0, 0, 0]$ --- the coefficient of
$u^4$ only, matching $B(4,4,u) = u^4$.
The lower-left triangle is all zeros because $B(4,i,u)$
contains $u^i$ as its lowest power.
\end{notebox}


% =============================================================
\needspace{8\baselineskip}
\section{$M \cdot G$ --- Polynomial Coefficient Matrix}
% =============================================================

\begin{minted}{text}
  M * G  - Polynomial Coefficient Matrix
  ----------------------------------------------------------------

         x    y  z
       ---  ---  -
  u^4   13  -28  0
  u^3  -32   64  0
  u^2   24  -48  0
  u^1    0   16  0
  u^0    0    0  1
\end{minted}

$M \cdot G$ is the $5 \times 3$ product of the basis matrix
and the geometry matrix.
Row $k$ gives the coefficient of $u^{4-k}$ for each axis.

\[
  M \cdot G = \begin{bmatrix}
    13  & -28 & 0 \\
   -32  &  64 & 0 \\
    24  & -48 & 0 \\
     0  &  16 & 0 \\
     0  &   0 & 1
  \end{bmatrix}
\]

Reading each column:
\begin{itemize}
  \item \textbf{$x$ column:} $[13, -32, 24, 0, 0]^T$
        --- coefficients of $13u^4 - 32u^3 + 24u^2$.
  \item \textbf{$y$ column:} $[-28, 64, -48, 16, 0]^T$
        --- coefficients of $-28u^4 + 64u^3 - 48u^2 + 16u$.
  \item \textbf{$z$ column:} $[0, 0, 0, 0, 1]^T$
        --- constant polynomial $z(u) = 1$.
\end{itemize}

\begin{examplebox}
\textbf{Verify the $x$ column by hand:}

$M \cdot G_x = M \cdot [0, 0, 4, 4, 5]^T$.
Multiplying row by row:

\[
  \text{Row 0 (}u^4\text{):}\quad
  1(0)+(-4)(0)+6(4)+(-4)(4)+1(5) = 0+0+24-16+5 = 13 \checkmark
\]
\[
  \text{Row 1 (}u^3\text{):}\quad
  (-4)(0)+12(0)+(-12)(4)+4(4)+0(5) = 0+0-48+16+0 = -32 \checkmark
\]
\[
  \text{Row 2 (}u^2\text{):}\quad
  6(0)+(-12)(0)+6(4)+0(4)+0(5) = 0+0+24+0+0 = 24 \checkmark
\]
\end{examplebox}


% =============================================================
\needspace{8\baselineskip}
\section{Resulting Polynomials}
% =============================================================

\begin{minted}{text}
  Resulting Polynomials
  ----------------------------------------------------------------

  x(u)  =  13u^4  - 32u^3  + 24u^2
  y(u)  =  -28u^4  + 64u^3  - 48u^2  + 16u
  z(u)  =  1
\end{minted}

These are the final monomial polynomials, read directly from
the columns of $M \cdot G$:

\[
  x(u) = 13u^4 - 32u^3 + 24u^2
\]
\[
  y(u) = -28u^4 + 64u^3 - 48u^2 + 16u
\]
\[
  z(u) = 1
\]

These are identical to the polynomials produced by all
previous scripts.
The three methods (algebraic expansion, symbolic CAS,
matrix method) converge to the same result.

\begin{notebox}
\textbf{Quick sanity checks:}

$x(0) = 0$, $x(1) = 13 - 32 + 24 = 5$.
The curve starts at $x=0$ and ends at $x=5 = P_4^x$. \checkmark

$y(0) = 0$, $y(1) = -28 + 64 - 48 + 16 = 4$.
The curve starts at $y=0$ and ends at $y=4 = P_4^y$. \checkmark

$z(u) = 1$ for all $u$.
The curve is entirely in the plane $z=1$. \checkmark
\end{notebox}


% =============================================================
\needspace{8\baselineskip}
\section{Evaluation Table}
% =============================================================

\begin{minted}{text}
  Evaluation Table  -  P(u) = U * (M * G)
  ----------------------------------------------------------------

  +=============================================================+
  | u       x(u)               y(u)               z(u)        |
  |=============================================================|
  | 0       0  (0.0000)        0  (0.0000)   1  (1.0000)      |
  | 0.2500  269//256  (1.0508) 121//64 (1.8906) 1  (1.0000)   |
  | 0.5000  45//16  (2.8125)   9//4  (2.2500)   1  (1.0000)   |
  | 0.7500  1053//256 (4.1133) 201//64 (3.1406) 1  (1.0000)   |
  | 1       5  (5.0000)        4  (4.0000)       1  (1.0000)  |
  +=============================================================+
\end{minted}

The evaluation table is identical in content to script~03's
Step~5 summary table.
Each cell shows the exact rational result followed by
the 4-decimal-place float in parentheses.

\begin{center}
\begin{tabular}{ccccc}
\toprule
$u$ & $x(u)$ exact & $x(u)$ decimal & $y(u)$ exact & $y(u)$ decimal \\
\midrule
$0$    & $0$              & 0.0000 & $0$       & 0.0000 \\
$0.25$ & $269/256$        & 1.0508 & $121/64$  & 1.8906 \\
$0.50$ & $45/16$          & 2.8125 & $9/4$     & 2.2500 \\
$0.75$ & $1053/256$       & 4.1133 & $201/64$  & 3.1406 \\
$1$    & $5$              & 5.0000 & $4$       & 4.0000 \\
\bottomrule
\end{tabular}
\end{center}

\begin{examplebox}
\textbf{Verify $x(0.5) = 45/16$:}

At $u = 0.5$, the $U$ vector is
$[\tfrac{1}{16}, \tfrac{1}{8}, \tfrac{1}{4}, \tfrac{1}{2}, 1]$.
The $x$ column of $M \cdot G$ is $[13, -32, 24, 0, 0]$.

\[
  \tfrac{1}{16}(13) + \tfrac{1}{8}(-32) + \tfrac{1}{4}(24)
  + \tfrac{1}{2}(0) + 1(0)
  = \tfrac{13}{16} - 4 + 6
  = \tfrac{13}{16} + 2
  = \tfrac{45}{16}
\]

$45 \div 16 = 2.8125$. \checkmark
\end{examplebox}


% =============================================================
\needspace{8\baselineskip}
\section{Summary: Script 03 vs Script 04}
% =============================================================

\textit{The same computation, two different reading experiences.}

\begin{center}
\begin{tabular}{lll}
\toprule
Aspect & Script~03 & Script~04 \\
\midrule
Purpose & Teach the method step by step & Show the result concisely \\
$M$ derivation & Shown column by column & Not shown \\
Dot products & Shown per $u$ value & Not shown \\
$G$ display & As a plain list & As a labelled matrix \\
Output length & Long --- every step visible & Short --- results only \\
Best used for & First reading, learning & Reference, verification \\
\bottomrule
\end{tabular}
\end{center}

\begin{notebox}
\textbf{Which should you use?}

Use script~03 when you want to understand \emph{why}
the numbers come out the way they do ---
when you are learning the matrix method or debugging.

Use script~04 when you already understand the method
and just want to \emph{see} the matrices and the final
curve polynomials quickly.

\textit{Script~03 is the lecture.
Script~04 is the cheat sheet.
Keep both.}
\end{notebox}

\vspace{2em}\hrule\vspace{0.5em}
{\small\textcolor{gray}{%
Output explanation for \texttt{04-bezier-matrix-straight.jl}, Case~2
\textbar{} Directed and curated by E.R.\ Nurwijayadi
\textbar{} Code and explanations AI-generated}}

\end{document}
